\SourceLeaf{245}%
as in the preceding article, provided that the symbols refer in general to an unspecified point on any of the shortest lines.

All the shortest lines of the same length \(r\) end on another line. We denote its length, measured from an arbitrary starting point, by \(v\). Thus \(v\) can be regarded as a function of the indeterminates \(r,\varphi\). If \(\lambda'\) denotes the point on the sphere corresponding to the direction of the element \(\mathrm{d}v\), and \(\xi',\eta',\zeta'\) denote its coordinates with respect to the center of the sphere, we have:
\[
\frac{\mathrm{d}x}{\mathrm{d}\varphi}=\xi'\mathord{.}\frac{\mathrm{d}v}{\mathrm{d}\varphi},\quad
\frac{\mathrm{d}y}{\mathrm{d}\varphi}=\eta'\mathord{.}\frac{\mathrm{d}v}{\mathrm{d}\varphi},\quad
\frac{\mathrm{d}z}{\mathrm{d}\varphi}=\zeta'\mathord{.}\frac{\mathrm{d}v}{\mathrm{d}\varphi}
\]
From this and from
\[
\frac{\mathrm{d}x}{\mathrm{d}r}=\xi,\quad\frac{\mathrm{d}y}{\mathrm{d}r}=\eta,\quad\frac{\mathrm{d}z}{\mathrm{d}r}=\zeta
\]
it follows that
\[
\begin{aligned}
\frac{\mathrm{d}x}{\mathrm{d}r}\mathord{.}\frac{\mathrm{d}x}{\mathrm{d}\varphi}+\frac{\mathrm{d}y}{\mathrm{d}r}\mathord{.}\frac{\mathrm{d}y}{\mathrm{d}\varphi}+\frac{\mathrm{d}z}{\mathrm{d}r}\mathord{.}\frac{\mathrm{d}z}{\mathrm{d}\varphi}
&=(\xi\xi'+\eta\eta'+\zeta\zeta')\mathord{.}\frac{\mathrm{d}v}{\mathrm{d}\varphi}\\
&=\cos\lambda\lambda'\mathord{.}\frac{\mathrm{d}v}{\mathrm{d}\varphi}
\end{aligned}
\]
We denote the left side of this equation, which is also a function of \(r,\varphi\), by \(S\). Differentiating it with respect to \(r\) gives:
\[
\begin{aligned}
\frac{\mathrm{d}S}{\mathrm{d}r}&=\frac{\mathrm{d}\mathrm{d}x}{\mathrm{d}r^2}\mathord{.}\frac{\mathrm{d}x}{\mathrm{d}\varphi}+\frac{\mathrm{d}\mathrm{d}y}{\mathrm{d}r^2}\mathord{.}\frac{\mathrm{d}y}{\mathrm{d}\varphi}+\frac{\mathrm{d}\mathrm{d}z}{\mathrm{d}r^2}\mathord{.}\frac{\mathrm{d}z}{\mathrm{d}\varphi}\\
&\quad+\tfrac{1}{2}\mathord{.}\frac{\mathrm{d}((\tfrac{\mathrm{d}x}{\mathrm{d}r})^2+(\tfrac{\mathrm{d}y}{\mathrm{d}r})^2+(\tfrac{\mathrm{d}z}{\mathrm{d}r})^2)}{\mathrm{d}\varphi}\\
&=\frac{\mathrm{d}\xi}{\mathrm{d}r}\mathord{.}\frac{\mathrm{d}x}{\mathrm{d}\varphi}+\frac{\mathrm{d}\eta}{\mathrm{d}r}\mathord{.}\frac{\mathrm{d}y}{\mathrm{d}\varphi}+\frac{\mathrm{d}\zeta}{\mathrm{d}r}\mathord{.}\frac{\mathrm{d}z}{\mathrm{d}\varphi}\\
&\quad+\tfrac{1}{2}\mathord{.}\frac{\mathrm{d}(\xi\xi+\eta\eta+\zeta\zeta)}{\mathrm{d}\varphi}
\end{aligned}
\]
But \(\xi\xi+\eta\eta+\zeta\zeta=1\), so the differential of this sum is \(=0\). If \(\rho\) again denotes the radius of curvature on the line \(r\), the preceding article gives
\[
\frac{\mathrm{d}\xi}{\mathrm{d}r}=\frac{X}{\rho},\quad\frac{\mathrm{d}\eta}{\mathrm{d}r}=\frac{Y}{\rho},\quad\frac{\mathrm{d}\zeta}{\mathrm{d}r}=\frac{Z}{\rho}
\]
We thus obtain
\[
\frac{\mathrm{d}S}{\mathrm{d}r}=\frac{1}{\rho}\mathord{.}(X\xi'+Y\eta'+Z\zeta')\mathord{.}\frac{\mathrm{d}v}{\mathrm{d}\varphi}=\frac{1}{\rho}\mathord{.}\cos L\lambda'\mathord{.}\frac{\mathrm{d}v}{\mathrm{d}\varphi}=0
\]
since \(\lambda'\) evidently lies on the great circle whose pole is \(L\). From this we conclude that \(S\) is independent of \(r\) and hence a function of \(\varphi\) alone. But for \(r=0\) we evidently have \(v=0\), and therefore also \(\frac{\mathrm{d}v}{\mathrm{d}\varphi}=0\) and \(S=0\), independently of \(\varphi\). It is therefore necessary that \(S=0\) in general, hence \(\cos\lambda\lambda'=0\), that is, \(\lambda\lambda'=90^0\). From this we conclude
